%%%%%%%%%%  scan idx 59  |  folio 46  %%%%%%%%%%
brique sur $X$ dont les sections sur un ouvert $V$ sont les sections de $F$ sur $U\cap V$. On constate facilement que $\overline F$ est \emph{quasi-cohérent} (i.e. défini sur tout ouvert affine $V$ -- donc ici sur $X$ tout entier -- par un module, de type fini ou non, sur l'anneau de coordonnées de l'ouvert en question). C'est immédiat si $U$ est lui-même affine, car $\overline F$ est alors défini sur l'ouvert $V$ par $\Gamma(U\cap V,F)$ considéré comme module sur $\Gamma(V,\cO_X)$. Dans le cas général, $U$ est réunion finie d'ouverts affines $U_i$, donc $\overline F$ est le noyau d'un homomorphisme évident de faisceaux quasi-cohérents~:
\[
 \sum_i \overline{F|U_i}\longrightarrow \sum_{i,j}\overline{F|U_i\cap U_j}
\]
donc est lui-même quasi-cohérent. Comme $X$ est supposé affine, $\overline F$ est donc engendré par ses sections. Comme $\overline F|U=F$ est cohérent, il existe un nombre fini de sections de $\overline F$ qui engendrent $\overline F|U$. Ces sections engendrent un sous-faisceau cohérent de $\overline F$, dont la restriction à $U$ est isomorphe à $F$.
\begin{flushright}cqfd.\end{flushright}

Supposons maintenant $X$ quasi-projective \emph{non singulière}. Alors, en vertu du théorème des syzygies, on a, avec les notations du corollaire au théorème 2.3, $\cF_0(X)=\cF(X)$. Dans ce cas, on a donc~:
\begin{equation}\tag{2.13}
 \underline{K}(\xi(X))=\underline{K}(\cF(X)).
\end{equation}
On peut exprimer ainsi cet isomorphisme~:

\textbf{Corollaire 2 au théorème 2.3.} \emph{Soit $X$ une variété quasi-projective non-singulière. Pour tout groupe abélien $A$, il y a correspondance biunivoque entre les fonctions additives $f$ de fibrés vectoriels sur $X$, à valeurs dans $A$, et les fonctions additives $g$ de faisceaux algébriques cohérents sur $X$, à valeurs dans $A$. À toute $f$ correspond l'unique $g$ telle que l'on ait~:}
